\section{Cronograma}
\label{sec:cronograma}

\subsection{Diagrama de Gantt}

\begin{figure}[H]
  \centering
  \begin{ganttchart}[
      hgrid,
      vgrid={*{3}{dotted}, black},
      title/.append style={fill=black, draw=black},
      title label font=\color{white}\small\bfseries,
      bar/.append style={fill=gray!40},
      bar label font=\small,
      group/.append style={fill=gray!70, draw=black},
      group label font=\small\bfseries,
      milestone/.append style={fill=black, draw=black},
      x unit=0.55cm,
      y unit title=0.6cm,
      y unit chart=0.55cm,
      bar height=0.5,
      group height=0.3
    ]{1}{16}
    %── Encabezado ─────────────────────────────────────────────
    \gantttitle{Semanas del semestre}{16} \\
    \gantttitlelist{1,...,16}{1} \\
    %── FASE 1 ─────────────────────────────────────────────────
    \ganttgroup{Fase 1: Recolección y Limpieza}{1}{5} \\
    \ganttbar{A1.1\ \ Descarga y consolidación SECOP II}{1}{3} \\
    \ganttbar{A1.2\ \ Limpieza y normalización de datos}{2}{4} \\
    \ganttbar{A1.3\ \ Construcción variable objetivo}{3}{5} \\
    %── FASE 2 ─────────────────────────────────────────────────
    \ganttgroup{Fase 2: Ingeniería de Características}{5}{8} \\
    \ganttbar{A2.1\ \ Selección y construcción de variables}{5}{7} \\
    \ganttbar{A2.2\ \ Pipeline de preprocesamiento}{6}{8} \\
    \ganttbar{A2.3\ \ Análisis exploratorio (EDA)}{6}{8} \\
    %── FASE 3 ─────────────────────────────────────────────────
    \ganttgroup{Fase 3: Entrenamiento y Evaluación}{8}{13} \\
    \ganttbar{A3.1\ \ Regresión Logística (línea base)}{8}{10} \\
    \ganttbar{A3.2\ \ Entrenamiento Random Forest}{9}{11} \\
    \ganttbar{A3.3\ \ Entrenamiento XGBoost}{10}{12} \\
    \ganttbar{A3.4\ \ Validación cruzada estratificada}{11}{13} \\
    %── FASE 4 ─────────────────────────────────────────────────
    \ganttgroup{Fase 4: Análisis y Cierre}{13}{16} \\
    \ganttbar{A4.1\ \ Comparación de métricas y modelos}{13}{14} \\
    \ganttbar{A4.2\ \ Análisis de importancia de variables}{13}{15} \\
    \ganttbar{A4.3\ \ Redacción monografía final}{14}{16} \\
    %── Hito ───────────────────────────────────────────────────
    \ganttmilestone{Entrega final}{16}
  \end{ganttchart}
  \caption{Diagrama de Gantt -- Cronograma del proyecto (16~semanas).}
  \label{fig:gantt}
\end{figure}

\subsection{Tabla de Hitos}

Los hitos funcionan como puntos de control internos. Si al llegar a uno el
entregable no está listo, hay margen para ajustar el ritmo antes de que el
retraso se propague a la fase siguiente.

\begin{table}[H]
  \centering
  \caption{Hitos del proyecto}
  \label{tab:hitos}
  \renewcommand{\arraystretch}{1.3}
  \begin{tabularx}{\linewidth}{c l X c}
    \toprule
    \textbf{N°} & \textbf{Hito} & \textbf{Criterio de cumplimiento} & \textbf{Semana} \\
    \midrule
    H1 & Datos listos
       & Dataset limpio, variable objetivo construida y umbral de \textit{proxy}
         documentado
       & 5 \\
    H2 & Pipeline funcional
       & Flujo de preprocesamiento serializado, probado y versionado en el
         repositorio
       & 8 \\
    H3 & Modelos entrenados
       & Los tres clasificadores ajustados con hiperparámetros registrados
         en el experimento
       & 12 \\
    H4 & Evaluación cerrada
       & Reporte de métricas con validación cruzada completo y revisado
       & 13 \\
    H5 & Modelo seleccionado
       & Elección justificada del algoritmo final con análisis SHAP incluido
       & 15 \\
    H6 & Entrega final
       & Monografía y repositorio de código públicos y documentados
       & 16 \\
    \bottomrule
  \end{tabularx}
\end{table}